\section{Analysis}
\label{sec:analysis}

\subsection{Why BOUNDARY Dominates}

The BOUNDARY strategy performs a deterministic sweep over a fixed set of
edge-case values: $\{0, 1, 2, -1, 10, 100\}$ and small primes, squares,
and Fibonacci numbers.
For number-theoretic conjectures of the form ``$P(n)$ holds for all $n$'',
false instances often fail at $n \in \{0, 1, 2\}$ by design (Tier~1) or by
nature (Tier~2 historical).
The 7B model's LLM-dependent strategies (SYMBOLIC\_PERTURBATION, ANALOGICAL)
generate large candidate sets (56--124 candidates per run, median ${\approx}100$) but none verify.

\paragraph{SymPy parser limitation.}
Post-hoc inspection of the raw loop results reveals that all 5 SP rounds in
E2 (and all ANALOGICAL rounds on the same conjectures) return
$\mathtt{holds=None}$ for every candidate; BOUNDARY also produces 0 verified
candidates for these conjectures.
The root cause is a parsing failure in \texttt{\_sympify\_eval\_number\_theory}:
the ``is divisible by'' regex captures the full NL preamble
(e.g.\ ``for all $n \geq 1$, $n^3-n$'') rather than just the expression,
causing a \texttt{SympifyError}; additionally, exponent notation
(\texttt{\^}) was not converted to Python's \texttt{**} before sympification.
Conjectures in the \emph{k divides f(n)} form (e.g.\ ``12 divides
$n(n{+}1)(n{+}2)$'') were unaffected and were correctly caught by BOUNDARY.
This bug has been patched in the codebase: a preamble-stripping regex now
extracts only the algebraic expression, and \texttt{\^} is normalized to
\texttt{**} before sympification.
The E1/E2 experimental results were obtained \emph{before} this patch.
E3 quantifies the recall impact: recall rises from $0.292$ (E1) to $0.389$ (E3),
a gain of $+7$ refutations.
All 7 new refutations are Tier~1 number-theory conjectures involving
$n^k$ expressions (e.g.\ ``$9$ divides $4^n-1$'', ``$25$ divides $n^5-n$'')
that previously returned \texttt{holds=None} due to the parsing failure.

\subsection{Strategist Early-Stop Bug and Fix}
\label{sec:strategist_fix}

During development we identified a bug in the Strategist's
\texttt{should\_stop} condition.
After processing several conjectures, global statistics for
RANDOM\_STRUCTURED, SYMBOLIC\_PERTURBATION, and ANALOGICAL accumulated
$\geq 5$ attempts with win rate $0.000$.
The low-confidence stopping condition
($\text{best win rate of untried strategies} < 0.05$)
then fired after a single BOUNDARY round for all subsequent conjectures.

\paragraph{Consequence.} Pre-patch: 69 of 71 survived conjectures received
only 1 round (BOUNDARY alone). Post-patch: all conjectures receive at least
2 distinct strategies before early stopping is permitted.

\paragraph{Impact on recall.} Zero.
Post-patch recall ($0.292$) is identical to pre-patch recall.
The additional RANDOM\_STRUCTURED rounds on 69 conjectures produced no new
refutations, confirming that the survived conjectures are genuinely hard for
the 7B model rather than merely under-explored.

\paragraph{Correct behavior.} The patch is nonetheless necessary: the correct
pipeline behavior is to try at least two strategies before concluding
low-confidence termination, and the pre-patch shortcut constituted incorrect
accounting of strategy coverage in experimental records.

\subsection{Float64 Precision Bug in RANDOM\_STRUCTURED}
\label{sec:evalf_bug}

Post-run analysis of E3 revealed that 2 of the 30 initially reported refutations
were false positives.
Both arose from the RANDOM\_STRUCTURED strategy at large candidate values
($n = 9{,}883$ and $n = 3{,}635$) on the conjectures
``$n^5 - n$ divisible by $10$'' and ``$n^5 - n$ divisible by $15$''.
The root cause is integer overflow via floating-point: the evaluator called
\texttt{int(expr.subs(n, n\_val).evalf())}, which converts the result to
\texttt{float64} before truncating to \texttt{int}.
For $n \gtrsim 3{,}000$, the value $n^5 - n$ exceeds $2^{53}$
(the float64 exact-integer range), causing rounding that corrupts the
last several decimal digits and produces a spurious non-zero remainder.

The fix removes \texttt{evalf()} from all integer evaluation paths
(\texttt{int(expr.subs(n, n\_val))} is exact for integer inputs);
the patched tests pass for all values including $n = 9{,}883$.
Neither affected conjecture should have been in the benchmark as false:
both are universally true by Fermat's little theorem ($n^5 \equiv n \pmod{10}$
and $n^5 \equiv n \pmod{15}$), constituting benchmark labeling errors.
E3 results in Table~\ref{tab:main_results} and Table~\ref{tab:breakdown}
exclude these two candidates; the reported precision $1.000$ reflects the
28 verified-correct refutations.

\subsection{Precision and False-Positive Rate}

Precision is $1.000$ in all runs ($\text{FPR} = 0.000$).
Every accepted counterexample (after excluding the two E3 false positives
described above) is SymPy-verified before being recorded.
SymPy integer arithmetic is exact for the polynomial and modular
expressions in these conjectures, making the zero false-positive rate
a structural guarantee rather than a statistical outcome.
The float64 bug described above does not affect BOUNDARY (which uses
$n \leq 100$, well within exact float64 range) or E1/E2 results.

\subsection{C-Agent Refinement Quality}

In E1, all 21 refuted conjectures have status \texttt{refined} in the raw
loop output (the original statement was refuted, then a refined version was
proposed and re-attacked); the evaluator correctly maps \texttt{refined} $\to$ refuted.
Inspection of several refined conjectures reveals a recurring failure mode:
the C-Agent proposes a refinement that still fails at $n=1$, and the loop
exhausts its refinement budget refuting successively more restricted statements
rather than converging on a true strengthening.
This points to a C-Agent prompt engineering opportunity.
E2 (32B) exhibits the same failure mode: boundary-exploitable conjectures
cycle through $n=1$ or $n=2$ counterexamples across all refinements, while
the C-Agent's proposed conditions are insufficiently tight. This is consistent
across model scales.
